\documentclass[10pt,a4paper]{amsart}
\usepackage[margin=19mm]{geometry}
\usepackage{amsmath,amssymb,mathtools}
\usepackage[scr=rsfs,cal=euler]{mathalfa}
\usepackage{xfrac}
\usepackage[unicode,hidelinks,bookmarks=false]{hyperref}
\newcommand{\ie}{i.e.\ }
\newcommand{\cf}{cf.\ }
\newcommand{\resp}{respectively\ }
\newcommand{\Subsection}{\subsection}
\providecommand{\nobreakdash}{\nobreak\hbox{-}}
\setlength{\emergencystretch}{2em}
\multlinegap0pt
\usepackage{xcolor}
\usepackage{amssymb}
\usepackage{mathtools}
\usepackage{tikz}
\newtheorem{lemma}{Lemma}
\newtheorem{theorem}{Theorem}
\newcommand{\DL}{\mathrm{DL}}
\newcommand{\Desc}{\mathrm{Desc}}
\newcommand{\Int}{\mathsf{Int}}
\newcommand{\Obs}{\mathsf{Obs}}
\newcommand{\Prob}{\mathbb{P}}
\newcommand{\calM}{\mathcal{M}}
\DeclareMathOperator{\doop}{do}
\title{The Causal Description Gap:\\ Information-Theoretic Separations Across Pearl's Hierarchy}
\author{Seyed Morteza Emadi}
\date{Source: arXiv:2605.02177v1 (CC BY 4.0)}
\begin{document}
\maketitle

\noindent\emph{Source and licence.} The Causal Description Gap:\\ Information-Theoretic Separations Across Pearl's Hierarchy. Version arXiv:2605.02177v1 (CC BY 4.0), distributed by arXiv under the Creative Commons Attribution 4.0 International licence, http://creativecommons.org/licenses/by/4.0/, as recorded in the arXiv licence metadata for this exact version and shown on the abstract page https://arxiv.org/abs/2605.02177v1. Full licence notice: \texttt{licenses/arxiv-2605.02177-CC-BY-4.0.txt}.\par\smallskip

\noindent\textit{Editorial scope.} Counted unit: the article's theorem thm:tree-separation (source0.tex lines 704--717). The article's complete original proof of it is delivered with it (source0.tex lines 719--740, one \texttt{proof} environment of 22 lines). No mathematics is written, repaired or re-derived: the statement and the proof are the article's own lines, in the article's own notation, and no retained line is substituted or re-flowed.  Retained uncounted as the source-local context the proof needs: lines 240--248; lines 272--275; lines 639--650; lines 663--673; lines 683--686.  Nothing was omitted from any retained span, so the package carries no omission record.  No retained line contains a citation, so the wrapper carries no bibliography and no bibliography entry.  No retained line contains a figure environment, an image-inclusion command or drawing code, so no image is delivered and the package declares no asset dependency.  Editorial formatting only, in addition: an AMS \texttt{amsart} preamble that restates the article's own declarations and macros used by the retained lines, namely the environments \texttt{lemma}, \texttt{theorem} on separate counters, together with the standard packages those lines need.  The retained environments print in this document as Lemma 1, Theorem 1 in order of appearance, whereas the article prints the same items with the numbering of its own full document.  The complete native source of the article as distributed in the arXiv e-print for this exact version is bundled unchanged as \texttt{sources/199703/source0.tex}.
\begin{lemma}[Counting bound for Kolmogorov complexity]
\label{lem:counting}
Let $z_1, \ldots, z_N$ be $N$ distinct strings and let $w$ be an arbitrary auxiliary string. Then:
\begin{enumerate}
    \item At least one $z_i$ satisfies $K(z_i \mid w) \geq \log_2 N - O(1)$.
    \item For every $c \geq 0$, all but at most a $2^{-c}$ fraction of the $z_i$ satisfy $K(z_i \mid w) \geq \log_2 N - c - O(1)$.
\end{enumerate}
The unconditional version ($w$ empty) is the special case. We apply this lemma with $w = n$, $w = (P, n)$, or $w = (I, n)$ as needed.
\end{lemma}

\begin{lemma}[Conditioning on a constant-complexity object is free]
\label{lem:conditioning-simple}
If $K(\Obs(M) \mid n) = O(1)$, then $K(\Int_1(M) \mid \Obs(M), n) = K(\Int_1(M) \mid n) \pm O(1)$, and consequently $\Delta_{2|1}(M) = \DL_2(M) \pm O(1)$.
\end{lemma}

\begin{lemma}[All Trees Have the Same Observations]
\label{lem:tree-obs}
For every rooted tree $T$, the observational distribution of $M_T$ is:
\[
P^\star(x_1, \ldots, x_n) = \begin{cases}
1/2 & \text{if } x_1 = x_2 = \cdots = x_n = 0, \\
1/2 & \text{if } x_1 = x_2 = \cdots = x_n = 1, \\
0 & \text{otherwise}.
\end{cases}
\]
Consequently, $\DL_1(M_T) = K(P^\star \mid n) = O(1)$.
\end{lemma}

\begin{lemma}[Interventions Reveal Descendants]
\label{lem:tree-intervention}
For any node $i \in [n]$ and the intervention $\doop(X_i = 0)$:
\[
P_{M_T}^{\doop(X_i = 0)}(X_j = 0) = \begin{cases}
1 & \text{if } j \in \Desc_T(i), \\
1/2 & \text{if } j \notin \Desc_T(i),
\end{cases}
\]
where $\Desc_T(i)$ is the set of descendants of $i$ in $T$ (including $i$ itself).
\end{lemma}

\begin{lemma}[Descendant Sets Determine the Tree]
\label{lem:descendants-determine-tree}
The collection of descendant sets $\{\Desc_T(i) : i \in [n]\}$ uniquely determines the rooted tree $T$.
\end{lemma}

\begin{theorem}[Rooted-Tree Separation: $\Theta(n \log n)$]
\label{thm:tree-separation}
Let $\calM_{\mathrm{tree}}^n := \{M_T : T \text{ is a rooted labeled tree on } [n]\}$. Then:
\begin{enumerate}
    \item The map $T \mapsto \Int_1(M_T)$ is injective on $\calM_{\mathrm{tree}}^n$.
    \item \textbf{Lower bound:} There exists $T$ with $\Delta_{2|1}(M_T) \geq (n-1)\log_2 n - O(\log n)$.
    \item \textbf{Upper bound:} For all $T$, $\DL_2(M_T) \leq (n-1)\log_2 n + O(\log n)$.
    \item \textbf{High probability:} For uniformly random rooted labeled tree $T$,
    \[
    \Prob\bigl[\Delta_{2|1}(M_T) \geq (n-1)\log_2 n - c - O(\log n)\bigr] \geq 1 - 2^{-c}.
    \]
\end{enumerate}
Consequently, $\Delta_{2|1}(M_T) = \Theta(n \log n)$ for all but an exponentially small fraction of trees.
\end{theorem}

\begin{proof}
\textbf{Part 1 (Injectivity):} By Lemma~\ref{lem:tree-intervention}, from $\Int_1(M_T)$ we can read off:
\[
\Desc_T(i) = \{j \in [n] : P_{M_T}^{\doop(X_i = 0)}(X_j = 0) = 1\}
\]
for each $i$. By Lemma~\ref{lem:descendants-determine-tree}, the descendant sets determine $T$. Thus different trees yield different interventional families.

\textbf{Part 2 (Lower bound):} By Cayley's formula, there are $n^{n-1}$ rooted labeled trees on $[n]$. (Cayley's formula gives $n^{n-2}$ labeled unrooted trees on $[n]$; choosing a root from $n$ vertices gives $n^{n-1}$ rooted labeled trees.) By injectivity, there are $n^{n-1}$ distinct elements in $\{\Int_1(M_T) : T \in \calM_{\mathrm{tree}}^n\}$.

By the counting bound for Kolmogorov complexity: among $N$ distinct strings, at least one has complexity $\geq \log_2 N - O(1)$. Applying this with $N = n^{n-1}$:
\[
\max_T K(\Int_1(M_T) \mid n) \geq \log_2(n^{n-1}) - O(1) = (n-1)\log_2 n - O(1).
\]

By Lemma~\ref{lem:tree-obs}, $\DL_1(M_T) = O(1)$ for all $T$. By Lemma~\ref{lem:conditioning-simple}, $\Delta_{2|1}(M_T) = \DL_2(M_T) \pm O(\log n)$.

\textbf{Part 3 (Upper bound):} Given tree $T$, we can encode it using a Pr\"ufer sequence: a sequence of $n-2$ labels from $[n]$, requiring $(n-2)\log_2 n$ bits, plus $O(\log n)$ bits to specify the root. Total: $(n-1)\log_2 n + O(\log n)$ bits.

A constant-length program can then simulate $M_T$ and output $\Int_1(M_T)$. Thus $K(\Int_1(M_T) \mid n) \leq (n-1)\log_2 n + O(\log n)$.

\textbf{Part 4 (High probability):} By part~(2) of the counting bound (Lemma~\ref{lem:counting}), among the $n^{n-1}$ distinct strings $\{\Int_1(M_T)\}$, all but a $2^{-c}$ fraction satisfy $K(\Int_1(M_T) \mid n) \geq (n-1)\log_2 n - c - O(1)$. Transferring to $\Delta_{2|1}$ via Lemma~\ref{lem:conditioning-simple} incurs an additional $O(\log n)$ term.
\end{proof}

\end{document}
